\BourbakiExerciseGroup

\begin{enumerate}
\def\labelenumi{\arabic{enumi})}
\tightlist
\item
  Let \(\mathbf{K}\) be a field of characteristic \(p > 0\) , \(\mathbf{E}\) an extension of \(\mathbf{K}\) and \(\mathbf{B}\) a \(p\) -basis of \(\mathbf{E}\) over \(\mathbf{K}(\mathbf{E}^p)\) (V, p.~98).
\end{enumerate}

\begin{enumerate}
\def\labelenumi{\alph{enumi})}
\item
  Assume that \(\mathbf{E} \subset \mathbf{K}^{p^{-1}}\) ; let \(\mathbf{K}_0\) be a subfield of \(\mathbf{K}\) such that \(\mathbf{E}\) is separable over \(\mathbf{K}_0\) . Show that the set \(\mathbf{B}^p\) is a \(p\) -independent subset over \(\mathbf{K}_0(\mathbf{K}^p)\) (observe that if \((a_{\lambda})\) is a basis of the vector \(K_{0}\) -space K, then \((a_{\lambda}^{p})\) is a basis of \(K_{0}(K^{p})\) over \(K_{0}\) . Let C be a subset of K, disjoint from \(B^{p}\) and such that \(B^{p} \cup C\) is a p-basis of K over \(K_{0}(K^{p})\) ; show that \(B \cup C\) is a p-basis of E over \(K_{0}(E^{p})\) .
\item
  Assume that \(E \subset K^{p^{-n}}\) and that E is separable over a subfield \(K_{0}\) of K. Show that if the degree of imperfection of K over \(K_{0}\) is finite (V, p.~170. Ex. 1), then it is equal to the degree of imperfection of E over \(K_{0}\) (use a)).
\end{enumerate}

\begin{enumerate}
\def\labelenumi{\arabic{enumi})}
\setcounter{enumi}{1}
\tightlist
\item
  Let \(\mathbf{K}\) be a field of characteristic \(p > 0\) , \(\mathbf{E}\) an extension of \(\mathbf{K}\) and \(\mathbf{F}\) a separable extension of \(\mathbf{E}\) . Show that any two of the following three properties imply the third:
\end{enumerate}

\begin{enumerate}
\def\labelenumi{\alph{enumi})}
\item
  B is a \(p\) -basis of E over \(\mathbf{K}(\mathbf{E}^p)\) ;
\item
  C is a \(p\) -basis of \(\mathbf{F}\) over \(\operatorname{E}(\mathbf{F}^p)\) ;
\item
  \(\mathbf{B} \subset \mathbf{E}\) , \(\mathbf{B} \cup \mathbf{C}\) is a \(p\) -basis of \(\mathbf{F}\) over \(\mathbf{K}(\mathbf{F}^p)\) and \(\mathbf{B} \cap \mathbf{C} = \emptyset\) . (Use the fact that if \((c_{\mu})\) is a basis of the vector E-space \(\mathbf{F}\) , then \((c_{\mu}^p)\) is a basis of \(\mathbf{K}(\mathbf{E}^p)[\mathbf{F}^p]\) over \(\mathbf{K}(\mathbf{E}^p)\) and of \(\mathbf{E}[\mathbf{F}^p]\) over \(\mathbf{E}\) .)
\end{enumerate}

\begin{enumerate}
\def\labelenumi{\arabic{enumi})}
\setcounter{enumi}{2}
\item
  Let K be a field of characteristic p \textgreater{} 0, E an extension of K, and F a finitely generated extension of E. Show that the degree of imperfection of F over K is at least equal to that of E over K (reduce to the following two cases: \(1^{\circ}\) F is separable over E; \(2^{\circ}\) F = E(x), where \(x^{p} \in E\) (cf.~V, p.~170, Ex. 2)).
\item
  Let \(\mathbf{F}\) be a separable extension of a field \(\mathbf{K}\) of characteristic \(p > 0\) . Show that if \(\mathbf{E}\) is a subextension of \(\mathbf{F}\) , relatively perfect over \(\mathbf{K}\) (V, p.~170, Ex. 1), \(\mathbf{F}\) is separable over \(\mathbf{E}\) .
\item
  Let \(\mathbf{K}\) be a field of characteristic \(p > 0\) . Show that if \(\mathbf{E}\) is an algebraic extension of \(\mathbf{K}\) or a relatively perfect extension of \(\mathbf{K}\) , then \(\mathrm{E}^{p^{-\infty}} = \mathrm{E}(\mathrm{K}^{p^{-\infty}})\) .
\item
  Let \(\mathbf{K}\) be a field of characteristic \(p > 0\) , \(\mathbf{E}\) a separable extension of \(\mathbf{K}\) and \(\mathbf{B}\) a \(p\) -basis of \(\mathbf{E}\) over \(\mathbf{K}(\mathbf{E}^p)\) .
\end{enumerate}

\begin{enumerate}
\def\labelenumi{\alph{enumi})}
\item
  Show that B is algebraically free over K (consider an algebraic relation of least degree between the elements of B and put the degrees of the variables which occur in the form \(kp + h\) with \(0 \leqslant h \leqslant p - 1\) ). Deduce that if tr. \(\deg_{K}E < +\infty\) , then the degree of imperfection of E over K is at most equal to tr. \(\deg_{K}E\) .
\item
  Show that E is separable and relatively perfect over K(B).
\end{enumerate}

\begin{enumerate}
\def\labelenumi{\arabic{enumi})}
\setcounter{enumi}{6}
\item
  Show that a relatively perfect transcendental extension E of a field K of characteristic p \textgreater{} 0 is not finitely generated over K. (In the opposite case show that for every transcendence basis B of E over K, E would be algebraic and separable over K(B).)
\item
  If E is a separable algebraic extension of a field K of characteristic p \textgreater{} 0, then the largest perfect subextension \(E^{p^{\infty}}\) of E (the intersection of the fields \(E^{p^{n}}\) for \(n \geqslant 0\) ) is algebraic over the largest perfect subfield \(K^{p^{\infty}}\) of K. (If \(x \in E^{p^{\infty}}\) , show that for every integer n \textgreater{} 0 we have \(x^{p^{-n}} \in \mathrm{K}(x)\) and deduce that the minimal polynomial of x over \(K^{p^{n}}\) is the same as its minimal polynomial over K.)
\item
  Let \(\mathbf{E}\) be a field, \(\mathbf{G}\) a group of automorphisms of \(\mathbf{E}\) and \(\mathbf{K}\) the subfield of \(\mathbf{E}\) consisting of the elements invariant under \(\mathbf{G}\) .
\end{enumerate}

\begin{enumerate}
\def\labelenumi{\alph{enumi})}
\item
  Show that for every subfield L of E which is stable for the elements of G, L and K are linearly disjoint over L ∩ K. (Consider a linear relation \(\sum_{j}\lambda_{j}x_{j}=0\) between the elements \(x_{j} \in \mathbf{K}\) which are linearly independent over \(\mathbf{L} \cap \mathbf{K}\) , with \(\lambda_{j} \in \mathbf{L}\) not all zero, the number of \(\lambda_{j} \neq 0\) being the least possible.)
\item
  If \(L \cap K\) is the field of invariants of a group of automorphisms of K, show that L is the field of invariants of a group of automorphisms of \(\mathrm{L}(\mathbf{K})\) .
\end{enumerate}

\begin{enumerate}
\def\labelenumi{\arabic{enumi})}
\setcounter{enumi}{9}
\item
  Let E be an extension of a field K, let G be a compact subgroup of the topological group \(\operatorname{Aut}_{\mathbf{K}}(\mathbf{E})\) and let F be the field of invariants of G. Show that E is a Galois extension of F and that G is canonically isomorphic to the topological group \(\operatorname{Aut}_{\mathbf{F}}(\mathbf{E})\) . (Observe that G operates in a continuous fashion on the discrete space E to establish that an element \(x \in E\) has a finite orbit under the action of G.)
\item
  Let \(\mathbf{K}\) be a field of characteristic \(p > 0\) . For every commutative K-algebra \(\mathbf{A}\) , raising to the \(p\) -th power makes the ring \(\mathbf{A}\) into an A-algebra, hence also a K-algebra which is denoted by \(\mathbf{A}^{p^{-1}}\) . Let \(\mathbf{A}\) be a commutative K-algebra. Show that the following conditions are equivalent:
\end{enumerate}

\begin{enumerate}
\def\labelenumi{(\roman{enumi})}
\item
  \(\mathbf{A}\) is a separable K-algebra.
\item
  The unique K-homomorphism \(\Phi: \mathbf{A} \otimes_{\mathbf{K}} \mathbf{K}^{p^{-1}} \to \mathbf{A}^{p^{-1}}\) such that \(\Phi(a \otimes \lambda) = a^p \cdot \lambda (a \in \mathbf{A}, \lambda \in \mathbf{K})\) is injective.
\item
  For every family \((b_{i})_{i\in I}\) of \(\mathbf{K}\) which is linearly free over \(\mathbf{K}^p\) and every family \((a_{i})_{i\in I}\) of \(\mathbf{A}\) the equality \(\sum_{i}a_i^p b_i = 0\) implies that \(a_{i} = 0\) for all \(i\) .
\end{enumerate}

(It may be verified that (ii) \(\Leftrightarrow\) (iii) and that (ii) is a consequence of V, p.~123, Th. 2, e).)

\begin{enumerate}
\def\labelenumi{\arabic{enumi})}
\setcounter{enumi}{11}
\tightlist
\item
  Let K be a field and \((\mathbf{X}_{i})_{i\in I}\) a family of indeterminates. Show that the algebra of formal power series \(\mathbf{K}[[(\mathbf{X}_{i})_{i\in I}]]\) is a separable K-algebra. (Apply Ex. 11; it may also be verified that for every finite extension \(K'\) of K, \(\mathbf{K}[[(\mathbf{X}_{i})_{i\in I}]]\otimes_{\mathbf{K}}K'\) is isomorphic to \(\mathbf{K}'[[(\mathbf{X}_{i})_{i\in I}]]\) and then Th. 2, d) of V, p.~123 applied.) Deduce that the field of fractions of \(\mathbf{K}[[(\mathbf{X}_{i})_{i\in I}]]\) is a separable extension of K (V, p.~120, Prop. 4).
\end{enumerate}

